\section{Light cones in the XY chain}\label{sec:xy}

\subsection{The chain}

Consider the spin chain on $\Z$ with Hamiltonian
\[
H_{XY}=\sum_n\bigl(X_nX_{n+1}+Y_nY_{n+1}\bigr)+\sum_nh_nZ_n,\qquad h_n=v(\theta+n\alpha),
\]
where $X_n,Y_n,Z_n$ are the Pauli matrices at site $n$. The interactions are nearest-neighbour and bounded, so the dynamics $\tau_t$ exists on the quasi-local algebra, and $\tau_t(A)=\lim_{L\to\infty}\tau^L_t(A)$ in norm for local $A$, where $\tau^L$ is generated by the restriction of $H_{XY}$ to $\Lambda_L=[-L,L]$~\cite{LiebRobinson1972}. All identities below are proved in finite volume and pass to the limit.

On $\Lambda_L$ the Jordan--Wigner transformation~\cite{LSM1961} $c_j=(\prod_{-L\le k<j}Z_k)\sigma^-_j$, $\sigma^-=(X-iY)/2$, gives $c_j^*c_j=(1+Z_j)/2$ and $H_{XY}|_{\Lambda_L}=\sum_{a,b}c_a^*(M_L)_{ab}c_b+\mathrm{const}$ with $M_L=-2(S+S^*)+2\operatorname{diag}(h)$ on $\ell^2(\Lambda_L)$. The diagonal unitary $W=\operatorname{diag}((-1)^j)$ gives $WM_LW=2H_L$, where $H_L$ is the Dirichlet restriction of $H=H_{v,\alpha,\theta}$. Hence
\[
\tau^L_t(c_x)=\sum_yU^L_t(x,y)c_y,\qquad|U^L_t(x,y)|=\bigl|\langle\delta_y,e^{-2itH_L}\delta_x\rangle\bigr| .
\]
For the Majorana operators $\gamma_{2x-1}=c_x+c_x^*$ and $\gamma_{2x}=i(c_x-c_x^*)$, $\tau_t(\gamma_{2x-1})$ and $\tau_t(\gamma_{2x})$ are real linear combinations of Majoranas whose total squared weight on the sites of any set $R$ is $\sum_{y\in R}|U_t(x,y)|^2$. Spin time $t$ therefore corresponds to one-body time $2t$, and we write
\[
p^{(x)}_s(R)=\sum_{y\in R}\bigl|\langle\delta_y,e^{-isH}\delta_x\rangle\bigr|^2
\]
for the one-body mass in $R$ at one-body time $s$. These formulas are read in the CAR algebra in finite volume; the even polynomials in Majoranas and the local spin observables used below have norm limits as $L\to\infty$.

\subsection{Exact commutator identities}

\begin{proposition}\label{prop:comm}
Let $R$ be a finite set of sites not containing $0$, $P_R=\prod_{k\in R}Z_k$ and $p=p^{(0)}_{2t}(R)$. Then
\[
\bigl\|[\tau_t(Z_0),P_R]\bigr\|=4\sqrt{p(1-p)} .
\]
In particular $\|[\tau_t(Z_0),Z_N]\|=4|u|\sqrt{1-|u|^2}$ with $u=\langle\delta_N,e^{-2itH}\delta_0\rangle$.
\end{proposition}

\begin{proof}
$\tau_t(n_0)=f^*f$ with $n_0=c_0^*c_0$ and $f=\sum_yU_t(0,y)c_y$, a normalized mode. Since $Z_k=2n_k-1=-(-1)^{n_k}$, $P_R=\pm(-1)^{N_R}$ with $N_R=\sum_{k\in R}n_k$, and conjugation by $P_R$ changes the sign of $c_y$ for $y\in R$ and fixes the others. Write $f=f_{\mathrm{in}}+f_R$ with $f_R$ the part supported in $R$. Then $P_RfP_R=f_{\mathrm{in}}-f_R$ and
\[
[f^*f,P_R]=2\bigl(f_{\mathrm{in}}^*f_R+f_R^*f_{\mathrm{in}}\bigr)P_R .
\]
Write $f_R=ag_1$ and $f_{\mathrm{in}}=bg_2$ with normalized modes $g_1,g_2$ of disjoint supports, $|a|^2=p$ and $|b|^2=1-p$. The operator $\bar ba\,g_2^*g_1+\bar ab\,g_1^*g_2$ vanishes on the zero- and two-particle sectors of $g_1,g_2$ and acts on the one-particle sector as $\bigl(\begin{smallmatrix}0&\bar ab\\a\bar b&0\end{smallmatrix}\bigr)$, so its norm is $|a||b|=\sqrt{p(1-p)}$. Since $P_R$ is unitary and $Z_0=2n_0-1$, the claim follows.
\end{proof}

\begin{proposition}[Infinite-temperature OTOC]\label{prop:otoc}
In finite volume $\Lambda\supset R\cup\{0\}$, with $C_R(t)=\frac12\,2^{-|\Lambda|}\tr\bigl(|[\tau_t(Z_0),P_R]|^2\bigr)$ and $p$ as in Proposition~\ref{prop:comm},
\[
C_R(t)=4p(1-p).
\]
\end{proposition}

\begin{proof}
From the proof of Proposition~\ref{prop:comm}, $|[\tau_t(Z_0),P_R]|^2=16\,|a|^2|b|^2\bigl(n_{g_1}(1-n_{g_2})+n_{g_2}(1-n_{g_1})\bigr)$, conjugated by $P_R$. In the tracial state, $g_1,g_2$ are orthonormal modes and $\langle n_{g_1}(1-n_{g_2})\rangle=1/4$.
\end{proof}

So an order-one commutator, or OTOC, between $Z_0$ and the parity of a far region is equivalent to a fixed fraction, bounded away from $0$ and $1$, of the one-body packet lying in that region.

\begin{remark}\label{rem:otoc}
The tracial OTOC for a single receiver site is contained in the general Gaussian-state formula of Riddell and S{\o}rensen~\cite[Appendix~C, eq.~(C33)]{RiddellSorensen2020} (numbering of the PDF of arXiv:1908.03292v3), whose tracial specialization is $1-4p(1-p)$. The reduction of XY dynamics to one-body dynamics is classical~\cite{LSM1961}; Hamza, Sims and Stolz used it to derive dynamical localization of disordered XY chains from one-body localization~\cite{HSS2012}.
\end{remark}

\subsection{From Abel averages to every time}

Let $K$ be a Jacobi operator on $\ell^2(\Z)$ with off-diagonal entries $1$ and real diagonal, $\psi_t=e^{-itK}\delta_x$, and $\mathfrak f(t)=\|\un_R\psi_t\|$ with $R=\{y:|y-x|>N\}$.

\begin{lemma}\label{lem:abel}
Let $N\ge1$. Then $\mathfrak f$ is $1$-Lipschitz, and for $T\ge1$ and $|t|\le T$,
\[
\mathfrak f(t)^2\le\bigl(4e^3T\,P(N,T)\bigr)^{2/3},\qquad P(N,T)=\frac2T\int_0^\infty e^{-2s/T}\mathfrak f(s)^2\,ds .
\]
\end{lemma}

\begin{proof}
$\frac d{dt}\mathfrak f^2=\langle\psi_t,i[K,\un_R]\psi_t\rangle$, and $[K,\un_R]$ lives on the two bonds crossing the boundary of $R$. Since $N\ge1$, the two inner boundary sites $x\pm N$ are distinct, and Cauchy--Schwarz gives $|(\mathfrak f^2)'|\le2\mathfrak f$, so $\mathfrak f$ is $1$-Lipschitz. If $\mathfrak f(t_0)=\sigma\in(0,1]$ with $0\le t_0\le T$, then $\mathfrak f\ge\sigma/2$ on $[t_0,t_0+\sigma/2]$, and $P(N,T)\ge\frac2Te^{-2(t_0+1/2)/T}\frac{\sigma^2}4\frac\sigma2\ge e^{-3}\sigma^3/(4T)$. Negative times follow from $e^{itK}\delta_x=\overline{e^{-itK}\delta_x}$, as $K$ is real.
\end{proof}

For $N=0$ the Lipschitz constant is $\sqrt2$ instead of $1$, which changes only the numerical constant. Splitting the Abel integral at $kT$ gives, for every $k>0$,
\begin{equation}\label{eq:abelsplit}
P(N,T)\le\sup_{0\le t\le kT}\mathfrak f(t)^2+e^{-2k}.
\end{equation}

\subsection{Locality for even observables}

An observable is \emph{even} if it commutes with the spin flip $\prod_kZ_k$ over its support.

\begin{lemma}[Locality]\label{lem:locality}
Let $A$ be even, supported in $[-l,l]$, with $\|A\|\le1$, and let $\|B\|\le1$.
\begin{enumerate}
\item Let $R$ be disjoint from $[-l,l]$ such that $R_-=R\cap(-\infty,-l)$ and $R_+=R\cap(l,\infty)$ are intervals, and let $B$ lie in the algebra generated by the even observables supported in $R_-$ and those supported in $R_+$; for instance a parity string $P_R$, or a product $B_-B_+$ of even factors. Then
\[
\|[\tau_t(A),B]\|\le2^{4l+4}(4l+2)\max_{|z|\le l}\sqrt{p^{(z)}_{2t}(R)} .
\]
\item Let $B$ be any observable supported in $[r,\infty)$ with $r>l$. Then the same bound holds with $R$ replaced by the half-line $[r,\infty)$. The mirror statement holds on the left.
\end{enumerate}
\end{lemma}

\begin{proof}
In case~2 use Jordan--Wigner strings from the right end of the volume; the one-body matrix changes only by diagonal signs, so the moduli above are unchanged. Then $B$, of any parity, lies in the CAR algebra generated by the modes of $[r,\infty)$, because the strings of its sites run to the right end inside $[r,\infty)$. In case~1 either representation works: an even spin observable supported in an interval is an even element of the CAR algebra of that interval, since its strings cancel in pairs inside the interval; so $B$ is an even element of the CAR algebra of $R$. In both cases $A$ is an even polynomial in the $2(2l+1)$ Majoranas of $[-l,l]$: $A=\sum_Sc_S\gamma_S$ over even subsets $S$, with at most $2^{4l+2}$ terms and $|c_S|\le\|A\|$. For $\mu\in S$ split $G_\mu=\tau_t(\gamma_\mu)=G^{\mathrm{in}}_\mu+G^{\mathrm{out}}_\mu$ into its parts on Majoranas outside and inside $R$ (respectively the half-line). These are real linear combinations, so $\|G^{\mathrm{in}}_\mu\|\le1$ and $\|G^{\mathrm{out}}_\mu\|=\sqrt{p^{(z)}_{2t}(R)}$ for the site $z$ of $\mu$. A product of an even number of linear elements supported outside $R$ commutes with every even element of the CAR algebra of $R$, and, in case~2, with every element. So $[\tau_t(\gamma_S),B]=[\prod_\mu G_\mu-\prod_\mu G^{\mathrm{in}}_\mu,B]$, and telescoping gives $\|\prod G_\mu-\prod G^{\mathrm{in}}_\mu\|\le\sum_\mu\|G^{\mathrm{out}}_\mu\|$. Sum over $S$.
\end{proof}

Case~2 needs the whole half-line: for an odd observable $B$, the Jordan--Wigner string of $B$ runs through every site between $B$ and the end of the chain, and a bound through the mass on the support of $B$ alone is not available. The bounds concern even source observables and the receiver classes above; they are not the usual exponential Lieb--Robinson bounds for arbitrary observables.

\subsection{Light-cone exponents}

For a site $x$ and $N\ge0$ let $\mathcal B(x,N)$ be the set of observables $B$ with $\|B\|\le1$ of the two kinds in Lemma~\ref{lem:locality}, relative to the site $x$: elements of the algebra generated by even observables of $R_-$ and of $R_+$, for sets $R\subset\{y:|y-x|>N\}$ with $R_\pm$ intervals on the two sides of $x$; and arbitrary observables supported in $[x+N+1,\infty)$ or in $(-\infty,x-N-1]$. Let $\mathcal A_l(x)$ be the set of even observables supported in $[x-l,x+l]$ with $\|A\|\le1$.

\begin{definition}\label{def:LR}
The \emph{light-cone exponent at the site $x$} of the chain with phase $\theta$ is
\[
\Gamma_x(\theta)=\inf\Bigl\{a\ge0:\ \forall l,m:\ \sup_{|t|\le T}\ \sup_{A\in\mathcal A_l(x),\,B\in\mathcal B(x,T^a)}\|[\tau_t(A),B]\|=O(T^{-m})\Bigr\},
\]
and the \emph{uniform light-cone exponent} $\Gamma_{\mathrm{unif}}(\theta)$ is defined in the same way with an additional supremum over $x\in\Z$ inside the $O(T^{-m})$.
\end{definition}

Translation by $x$ sites maps the chain with phase $\theta$ to the chain with phase $\theta+x\alpha$, so $\Gamma_x(\theta)=\Gamma_0(\theta+x\alpha)$. By the ballistic bound $|\langle\delta_y,e^{-isH}\delta_x\rangle|\le(e|s|)^{|y-x|}/|y-x|!$, all these exponents are at most $1$.

\begin{theorem}[The fixed-site light cone]\label{thm:LRfixed}
Let $\lambda>1$, $\alpha$ irrational, and $h_n=2\lambda\cos2\pi(\theta+n\alpha)$. Let $\Ga^\Z=\bigcap_{x\in\Z}(\Ga-x\alpha)$, a set of full measure, and $\gamma=\max\{0,1-L/\beta(\alpha)\}$.
\begin{enumerate}
\item Let $\beta(\alpha)<\infty$ and $\theta\in\Ga^\Z$. For every $a>\gamma$, $l\ge0$ and $m>0$ there is $C$ with $\|[\tau_t(A),B]\|\le CT^{-m}$ for all $T\ge1$, $|t|\le T$, $A\in\mathcal A_l(0)$ and $B\in\mathcal B(0,T^a)$.
\item Let $\beta(\alpha)>L$ and $\theta$ arbitrary. For every $a<\gamma$ there are times $t_j\to\infty$ and finite sets $R_j\subset\{y:|y|>t_j^a\}$, each the union of two intervals, with
\[
\bigl\|[\tau_{t_j}(Z_0),P_{R_j}]\bigr\|\ge\sqrt{2c_B}=2^{-15.5}(1+4\pi\lambda)^{-1/2},\qquad c_B=\frac{2^{-32}}{1+4\pi\lambda}.
\]
\end{enumerate}
Consequently $\Gamma_x(\theta)=\max\{0,1-L/\beta(\alpha)\}$ for every site $x$ and every $\theta\in\Ga^\Z$, read as $1$ when $\beta(\alpha)=\infty$.
\end{theorem}

\begin{proof}
1. Fix $a'\in(\gamma,a)$. For $|z|\le l$ the phase $\theta+z\alpha$ lies in $\Ga$, and~\cite[Theorem~A]{DMamo} gives $P_{\theta+z\alpha}(S^{a'},S)=O(S^{-m'})$ for every $m'$. For large $T$, $T^a-l\ge(2T)^{a'}$. Lemma~\ref{lem:abel} at time scale $2T$ gives $\sup_{|s|\le2T}p^{(z)}_s(\{y:|y-z|>(2T)^{a'}\})=O(T^{-m''})$ for every $m''$. Lemma~\ref{lem:locality}, with finitely many $z$, finishes the proof; the constants depend on $\theta$ through the indices from which the sine conditions hold.

2. Choose $L<h<c<\beta$ with $1-h/c>a$, and scales $n_j$ with $b_{n_j}\ge c$ and $b_{n_j}\to\beta$. At $T_j=|\delta_{n_j}|^{-1}$ and $N_j=T_je^{-hq_{n_j}}=T_j^{1-h/b_{n_j}}$, \cite[Theorem~B]{DMamo} gives $P(N_j,T_j)\ge c_B$ at every phase. By~\eqref{eq:abelsplit} with $e^{-2k}=c_B/2$, some one-body time $s\le kT_j$ has $p^{(0)}_s(\{|y|>N_j\})\ge c_B/2$. Since $p_0=0$ and $s\mapsto p_s$ is continuous, there is $s_j\le s$ with $p^{(0)}_{s_j}(\{|y|>N_j\})=c_B/2$. Truncating to $R_j=\{N_j<|y|\le M_j\}$ with $M_j$ large keeps the mass in $[c_B/4,c_B/2]$, and then $4\sqrt{p(1-p)}\ge\sqrt{2c_B}$. Proposition~\ref{prop:comm} at spin time $t_j=s_j/2$ gives the bound. The times tend to infinity, because a fixed fraction of the packet cannot reach distance $N_j\to\infty$ in bounded time, by the ballistic bound. Finally $N_j\ge(kT_j)^a\ge t_j^a$ for large $j$, since $1-h/b_{n_j}>a$ eventually.

For the last statement, $\Gamma_x(\theta)=\Gamma_0(\theta+x\alpha)$ and $\theta+x\alpha\in\Ga^\Z$. Part~1 gives $\Gamma_0\le\gamma$ when $\beta<\infty$, and part~2 gives $\Gamma_0\ge\gamma$ when $\beta>L$; when $\beta\le L$, $\gamma=0$. When $\beta=\infty$, part~2 applies with any $a<1$.
\end{proof}

\begin{remark}\label{rem:LRprior}
Damanik, Lemm, Lukic and Yessen identified anomalous Lieb--Robinson exponents of the Fibonacci XY chain with upper one-body transport exponents and treated uniformity over the hull~\cite{DLLY2016}. Kachkovskiy proved velocity lower bounds for isotropic quasi-periodic XY chains in an absolutely continuous regime~\cite{Kachkovskiy2016}. Theorem~\ref{thm:LRfixed} is the corresponding statement for the Aubry--Andr\'e chain with Liouville frequencies, with its exact arithmetic exponent.
\end{remark}

The next lemma turns a polynomially small one-body tail into a commutator bound.

\begin{lemma}\label{lem:tailcomm}
Let $\varphi_n$ be phases, $T_n\to\infty$ and $N_n\ge T_n^{a_*-o(1)}$ with $a_*>0$, and suppose $P_{\varphi_n}(N_n,T_n)\ge T_n^{-C}$. Then for every $a<a_*$ there are $C'$, times $t_n\to\infty$ and finite sets $R_n\subset\{y:|y|>t_n^a\}$, each a union of two intervals, such that the chain with phase $\varphi_n$ satisfies $\|[\tau_{t_n}(Z_0),P_{R_n}]\|\ge t_n^{-C'}$.
\end{lemma}

The phases may vary with $n$, and a constant sequence is allowed; the proof treats each $n$ separately, and its only other input, the ballistic bound, holds uniformly in the phase because $\|v\|_\infty$ bounds the field at every phase.

\begin{proof}
By~\eqref{eq:abelsplit} with $e^{-2k}=T_n^{-C}/2$, some one-body time $s\le CT_n\log T_n$ has $p^{(0)}_s(\{|y|>N_n\})\ge T_n^{-C}/2$. Take the first time $s_n$ at which this mass equals $T_n^{-C}/4$, and truncate to a finite annulus $R_n=\{N_n<|y|\le M\}$ with mass $p\in[T_n^{-C}/8,T_n^{-C}/4]$. Proposition~\ref{prop:comm} gives $\|[\tau_{s_n/2}(Z_0),P_{R_n}]\|\ge4\sqrt{p/2}\ge T_n^{-C/2}$. The ballistic bound forces $s_n\ge cN_n$, so $T_n^{-C}$ is bounded below by a negative power of $t_n=s_n/2$; and $t_n\le CT_n\log T_n$, so $N_n\ge t_n^a$ for large $n$.
\end{proof}

\begin{theorem}[The uniform light cone]\label{thm:LRunif}
\begin{enumerate}
\item Let $v$ be real analytic with $H_{v,\alpha}$ supercritical, $\beta(\alpha)<\infty$, and $h_n=v(\theta+n\alpha)$. Then $\Gamma_{\mathrm{unif}}(\theta)\le\FU(\beta(\alpha))$ for every $\theta$.
\item Let $h_n=2\lambda\cos2\pi(\theta+n\alpha)$ with $\lambda>1$. Then $\Gamma_{\mathrm{unif}}(\theta)\ge\max\{0,1-L/(2\beta_e(\alpha))\}$ for every $\theta$, read as $1$ if $\beta_e=\infty$.
\item If moreover $\beta_e(\alpha)=\beta(\alpha)\in[L,\infty)$, then for every $\theta$
\[
\Gamma_{\mathrm{unif}}(\theta)=1-\frac L{2\beta(\alpha)},
\]
while $\Gamma_x(\theta)=1-L/\beta(\alpha)$ at every site $x$ when $\theta\in\Ga^\Z$.
\end{enumerate}
\end{theorem}

\begin{proof}
1. Let $a>a'>\FU(\beta)$. By Theorem~\ref{thm:U}, $\sup_\varphi P_\varphi((2T)^{a'},2T)=O(T^{-m})$ for every $m$. Lemma~\ref{lem:abel} turns this into $\sup_z\sup_{|s|\le2T}p^{(z)}_s(\{y:|y-z|>(2T)^{a'}\})=O(T^{-m'})$, uniformly in the site $z$, since its inequality is pointwise. For $A\in\mathcal A_l(x)$ and $B\in\mathcal B(x,T^a)$, the relevant sets lie in $\{y:|y-z|>T^a-l\}$ for $|z-x|\le l$, and Lemma~\ref{lem:locality} gives the bound uniformly in $x$.

2. Let $\beta_e>L/2$ and $h>L$ with $h/2+(h-L)<\beta_e$. For an even scale $n$ with $b_n\ge h/2+(h-L)$, the $q_{n+1}$ points $\theta+x\alpha$, $0\le x<q_{n+1}$, have gaps at most $\|q_n\alpha\|+\|q_{n+1}\alpha\|\le2/q_{n+1}$, and $1/q_{n+1}<2|\delta_n|$; so some $x_n$ with $0\le x_n<q_{n+1}<T_n=|\delta_n|^{-1}$ has $\|\theta+x_n\alpha-1/4\|\le2|\delta_n|$. Theorem~\ref{thm:quarter} at the phase $\theta+x_n\alpha$ gives $P_{\theta+x_n\alpha}(N_n,T_n)\ge T_n^{-2}e^{-2hq_n}$ with $N_n=\lfloor c_\eta T_ne^{-hq_n/2}\rfloor$. Lemma~\ref{lem:tailcomm}, applied to the chain at phase $\theta+x_n\alpha$, which is the chain at phase $\theta$ seen from the site $x_n$, gives commutators $\|[\tau_{t_n}(Z_{x_n}),P_{R_n}]\|\ge t_n^{-C'}$ with $R_n\subset\{y:|y-x_n|>t_n^a\}$ for every $a<1-h/(2\beta_e)$, along even scales with $b_n\to\beta_e$. So $\Gamma_{\mathrm{unif}}(\theta)\ge1-h/(2\beta_e)$; let $h\downarrow L$. If $\beta_e=\infty$ the same argument gives $1$.

3. Parts~1 and~2, with $\FU(\beta)=1-L/(2\beta)$ for the almost Mathieu operator when $\beta\ge L$ (Corollary~\ref{cor:typeI}), and Theorem~\ref{thm:LRfixed}.
\end{proof}

So uniformity over sites changes the light-cone exponent: when $\beta_e=\beta\in[L,\infty)$, almost every phase has fixed-site exponent $1-L/\beta$ at every site, and uniform exponent $1-L/(2\beta)$. At $\beta=\infty$ both equal $1$.

\subsection{Uncomputable light cones}

A real $x$ is \emph{right-c.e.\ relative to $\mathbf 0'$} if $\{r\in\Q:r>x\}$ is computably enumerable relative to the halting problem. Let $K_2$ be the halting set relative to $\mathbf 0'$ and $s_*=2-\sum_{e\in K_2}4^{-(e+2)}\in[23/12,2]$, a real of Turing degree $\mathbf 0''$~\cite[proof of Theorem~D]{DMamo}.

\begin{corollary}[A computable chain with a light cone of degree $\mathbf 0''$]\label{cor:LRdeg}
Let $\alpha_*$ be the explicit computable irrational of~\cite[Theorem~D(3)]{DMamo}, whose partial quotients after the first are even, so that all its denominators are odd, and $\beta(\alpha_*)=s_*\log2$. Consider the chain with $h_n=4\cos2\pi(1/4+n\alpha_*)=-4\sin(2\pi n\alpha_*)$. At every site $x$, the light-cone exponent $\Gamma_x(1/4)=1-1/s_*\in[11/23,1/2]$ has Turing degree $\mathbf 0''$. As $\alpha$ ranges over the computable irrationals with eventually odd denominators, the exponents $\Gamma_0(1/4)$ of the chains $h_n=4\cos2\pi(1/4+n\alpha)$ range over exactly the reals in $[0,1]$ that are right-c.e.\ relative to $\mathbf 0'$. For every computable $\lambda'$ with $1<\lambda'<2^{23/12}$, the chain $h_n=2\lambda'\cos2\pi(1/4+n\alpha_*)$ has $\Gamma_x(1/4)=1-\log\lambda'/\beta(\alpha_*)$ at every site, again of degree $\mathbf 0''$.
\end{corollary}

\begin{proof}
For odd $q_n$ and every fixed $x$, $|\sin(2\pi q_n(1/4+x\alpha))|=|\cos(2\pi x\delta_n)|\ge1/2$ for large $n$. So $1/4\in\Ga^\Z$ whenever the denominators are eventually odd, and Theorem~\ref{thm:LRfixed} identifies $\Gamma_x(1/4)$ with the one-body exponent of~\cite[Theorem~D]{DMamo}, including the value $1$ when $\beta=\infty$. The degree and value-set statements are~\cite[Theorem~D(2),(3)]{DMamo}. For $\lambda'$: since $\beta(\alpha_*)\ge(23/12)\log2>\log\lambda'$, the exponent $x'=1-\log\lambda'/\beta$ is positive, and $\beta=\log\lambda'/(1-x')$; both maps are computable for computable $\lambda'$, so $x'$ and $\beta$ have the same degree.
\end{proof}

The value-set statement concerns this family, with $\lambda=2$, phase $1/4$ and eventually odd denominators; it is not a classification over all computable fields and phases.

\begin{corollary}[Two exponents of degree $\mathbf 0''$ in one computable chain]\label{cor:twoexp}
Let $(v_m)$ be a computable sequence of rationals in $[1,2]$ with $\limsup v_m=s_*$~\cite[Lemma~10.1]{DMamo}. Define $\alpha=[0;a_1,a_2,\dots]$ by $a_1=2$ and, with $n_1=1$ and $n_{m+1}=n_m+3$,
\[
a_{n_m+1}=2^{\lceil v_mq_{n_m}\rceil},\qquad a_{n_m+2}=a_{n_m+3}=1 .
\]
Then $\alpha$ is computable, $\beta_e(\alpha)=\beta(\alpha)=s_*\log2$ and $\beta_o(\alpha)=0$. For the chains $h_n=4\cos2\pi(\theta+n\alpha)$:
\begin{enumerate}
\item for every $\theta$, $\Gamma_{\mathrm{unif}}(\theta)=1-1/(2s_*)\in[17/23,3/4]$;
\item for every $\theta\in\Ga^\Z$, in particular for almost every $\theta$, $\Gamma_x(\theta)=1-1/s_*\in[11/23,1/2]$ at every site $x$;
\item at $\theta=1/4$, $\Gamma_x(1/4)=1-1/(2s_*)$ at every site $x$;
\item there is a computable phase $\theta_*\in\Ga^\Z$.
\end{enumerate}
Both exponent values, $1-1/(2s_*)$ and $1-1/s_*$, have Turing degree $\mathbf 0''$. In particular the computable chain with phase $\theta_*$ has uniform exponent $1-1/(2s_*)$ and fixed-site exponent $1-1/s_*$ at every site: two different exponents of degree $\mathbf 0''$.
\end{corollary}

\begin{proof}
\emph{Arithmetic.} $q_{n_1}=q_1=2$ is even. If $q_{n_m}$ is even, then $q_{n_m-1}$ is odd, $q_{n_m+1}$ and $q_{n_m+2}$ are odd, and $q_{n_m+3}=q_{n_m+2}+q_{n_m+1}$ is even. The large partial quotient gives $b_{n_m}=v_m\log2+O(\log q_{n_m}/q_{n_m})$, while $b_n\to0$ along all other $n$. So $\beta=\beta_e=s_*\log2\ge\log2=L$ and $\beta_o=0$. The partial quotients are computable, so $\alpha$ is computable, although $\limsup v_m$ is not.

Part~1 is Theorem~\ref{thm:LRunif}(3), and part~2 is Theorem~\ref{thm:LRfixed}.

Part~3. Fix $x$. The first block to the right of $x$ in the chain with phase $1/4$ has phases $1/4+(x+1+k)\alpha$, and the rational block with phases $1/4+(x+1+k)p_n/q_n$ is a turning block, since $\cos(2\pi q_n(1/4+(x+1)p_n/q_n))=(-1)^{q_n/2}$ for even $q_n$. The site discrepancies are $(x+1+k)\delta_n/q_n$, at most $2|\delta_n|$ once $q_n>|x|$, so the proof of Theorem~\ref{thm:quarter} applies with the same exponents and gives $\au(1/4+x\alpha)\ge1-L/(2\beta_e)$. Lemma~\ref{lem:tailcomm} gives $\Gamma_x(1/4)\ge1-1/(2s_*)$, and part~1 gives the reverse inequality.

Part~4. Let $E_n=\{\theta:|\sin2\pi q_n\theta|<q_n^{-1/2}\}$, a finite union of intervals with computable endpoints and $|E_n|\le q_n^{-1/2}$. Since $q_n$ grows at least like the Fibonacci numbers, there is an explicit $N$ with $\sum_{n\ge N}|E_n|<1/2$, with an effective tail bound. Hence the measure of $(\bigcup_{n\ge N}E_n)\cap I$ is computable for every dyadic interval $I$, by finite unions and the tail bound. Start from $I_0=\T$; given $I_k$ whose intersection with the closed set $F=\T\setminus\bigcup_{n\ge N}E_n$ has positive measure, compute the measures for the two dyadic children until one is certified positive, which must happen, and let $I_{k+1}$ be that child. The nested intervals define a computable $\theta_*$. If $\theta_*$ were in the open set $\bigcup_{n\ge N}E_n$, a neighbourhood of $\theta_*$, hence some $I_k$, would lie in it, contradicting $|I_k\cap F|>0$. So $|\sin2\pi q_n\theta_*|\ge q_n^{-1/2}$ for $n\ge N$. For each fixed $x$, $q_n(\theta_*+x\alpha)=q_n\theta_*+xp_n+x\delta_n$, and the sine changes by at most $2\pi|x||\delta_n|\le2\pi|x|/q_n$, which is eventually below $q_n^{-1/2}/2$. So $\theta_*+x\alpha\in\Ga$ for every $x$.

\emph{Degrees.} $1-1/s_*$ and $1-1/(2s_*)$ are computable functions of $s_*$ on $[23/12,2]$ with computable inverses, so both have the degree of $s_*$, which is $\mathbf 0''$.
\end{proof}

Counting lattice sites is not averaging over phases: at the phase $1/4$ every site sees the exceptional exponent, while at almost every phase every site sees the typical one.
